\begin{table}[htbp]
\centering
\caption{Model-implied dispensary-revenue benchmark, Colorado border counties, 2024}
\label{tab:counterfactual}
\footnotesize
\begin{tabular}{lc}
\toprule
Quantity & Value \\
\midrule
Dispensary-density interaction coefficient $\theta$ & $-0.932$ \\
Free-entry benchmark multiplier $e^{-\theta}$ & 2.54 \\
\midrule
Actual 2024 revenue, CO border counties & \$324.2 M \\
Free-entry-benchmark counterfactual revenue (NM) & \$823.0 M \\
Model-implied additional dispensary revenue & +\$498.8 M \\
95\% county-cluster bootstrap interval & (\$-1.5 M, +\$2{,}526.0 M) \\
\bottomrule
\end{tabular}
\begin{minipage}{0.98\textwidth}
\vspace{0.5em}

\fontsize{9pt}{11pt}\selectfont
\textit{Notes:} The benchmark uses the \emph{dispensary-density} interaction
($\theta=-0.932$, Table~\ref{tab:regime_interaction}), the outcome for which the
difference across entry regimes is statistically significant. It asks how much
more dispensary revenue Colorado's border counties would have generated in 2024
if their density border premium had matched New Mexico's rather than the
restricted-entry premium, holding revenue per store at its observed Colorado
level. The multiplier $e^{-\theta}$ is applied to each Colorado border county's
own 2024 dispensary revenue and summed; New Mexico is the free-entry benchmark
state and contributes only to the estimation of $\theta$. Colorado is restricted
to 2022Q2 onward (Section~\ref{sec:crossstate}). The bootstrap interval resamples
counties with replacement within each of the two states (500 replicates),
refitting the density regression each time; it is wide and reaches down to zero
because $\theta$ is only marginally significant ($p\approx0.06$). The dollar
figure should therefore be read as an illustrative benchmark rather than an
identified estimate of foregone revenue.
\end{minipage}
\end{table}
